\documentclass[../Main.tex]{subfiles}

\begin{document}

\section{Design und Umsetzung des Kampfabschnitts} \label{sec:Umsetzung}

In diesem Kapitel werden zunächst Designvorschläge für die beiden inhaltlichen Komponenten der Erkenntnisfrage (Abschnitt~\ref{sec:Erkenntnisfrage}) aufgestellt: Gegnerverhalten und Kampfgefühl einerseits, Gegnerplatzierung und Leveldesign andererseits. Sie zielen auf die drei dort genannten spielerischen Qualitäten ab, nämlich einen nachvollziehbaren Spielfluss, eine angemessene Schwierigkeit und ein befriedigendes Spielgefühl. Dazu werden die Erkenntnisse aus der Referenzspielanalyse (Kapitel~\ref{sec:Referenzspielanalyse}) und der Leveldesign-Recherche (Kapitel~\ref{sec:Leveldesign}) zu spezifischen Designvorschlägen für unseren Kampfabschnitt verdichtet. Anschliessend wird für jeden Vorschlag beschrieben, wie er im Prototyp umgesetzt wurde und wo sich die Umsetzung überprüfen lässt. Die Überprüfung mittels Playtests folgt in Kapitel~\ref{sec:Evaluation}.

\subsection{Von der Recherche zu den Designvorschlägen} \label{sec:Designvorschlaege}

Jeder Designvorschlag folgt demselben Schema: der Vorschlag selbst, sein Nutzen, also welcher Aspekt der Erkenntnisfrage damit bedient wird, und die Herleitung aus Recherche und Spielanalysen. Die Vorschläge sind mit V1 bis V8 durchnummeriert. Dieselben Kürzel stehen hinter den Aussagen des Post-Questionnaire (Abschnitt~\ref{sec:PostQuestionnaire}) und den Kriterien des Beobachtungsrasters (Abschnitt~\ref{sec:Beobachtungsraster}), sodass jeder Vorschlag gezielt überprüft werden kann. Tabelle~\ref{tab:Designvorschlaege} gibt einen Überblick.

\begin{table}[H]
    \centering
    \small
    \renewcommand{\arraystretch}{1.3}
    \begin{tabular}{|>{\raggedright\arraybackslash}p{0.6cm}|>{\raggedright\arraybackslash}p{4.3cm}|>{\raggedright\arraybackslash}p{6.6cm}|>{\raggedright\arraybackslash}p{1.8cm}|}
        \hline
        \textbf{Nr.} & \textbf{Designvorschlag} & \textbf{Herleitung} & \textbf{Playtest} \\
        \hline
        \multicolumn{4}{|l|}{\textit{Gegnerverhalten (Abschnitt~\ref{sec:UmsetzungGegner})}} \\
        \hline
        V1 & Überlegenheit als Grundzustand & Kontrast \game{NieR: Automata} / \game{Ninja Gaiden} (Abschnitte~\ref{sec:ReferenzspielanalyseNier} und~\ref{sec:ReferenzspielanalyseNinja}); emotionsgeleitetes Design (Abschnitt~\ref{sec:Pacing}) & Q1, B1 \\
        \hline
        V2 & Begrenzung gleichzeitiger Angreifer & Beide Referenzspiele (Abschnitt~\ref{sec:Vergleich}); Pacing (Abschnitt~\ref{sec:Pacing}) & Q2, B2 \\
        \hline
        V3 & Taktisches Dilemma durch zwei Gegnerrollen & Frage D2 beider Analysen; Gegnerrollen im Encounter-Design (Abschnitt~\ref{sec:Encounter}) & Q3, B3 \\
        \hline
        V4 & Klar telegrafierte Angriffe & \game{NieR: Automata}, Frage M2; Lesbarkeit (Abschnitt~\ref{sec:Encounter}) & Q4, B4 \\
        \hline
        \multicolumn{4}{|l|}{\textit{Gegnerplatzierung und Kampfraum (Abschnitt~\ref{sec:UmsetzungLevel})}} \\
        \hline
        V5 & Lesbare Platzierung, dynamisches Verhalten & \game{NieR: Automata}, Frage A3; \game{Ninja Gaiden}, Frage D3 (Abschnitt~\ref{sec:Vergleich}); Lesbarkeit (Abschnitt~\ref{sec:Encounter}) & Q5, B6 \\
        \hline
        V6 & Der Raum als Druckmittel & \game{Ninja Gaiden}, Frage M3; räumliche Typologien (Abschnitt~\ref{sec:Raumtypologien}) & Q6, B7 \\
        \hline
        V7 & Wellenstruktur mit Erholungsphasen & \game{Ninja Gaiden}, Frage D1; \game{NieR: Automata}, Frage A4; Rollercoaster-Pacing (Abschnitt~\ref{sec:Pacing}) & Q7, B8 \\
        \hline
        V8 & Gestaffelte Einführung neuer Gegner & Scaffolding (Abschnitt~\ref{sec:Scaffolding}); \game{NieR: Automata}, Frage D1 & Q8, B9 \\
        \hline
    \end{tabular}
    \caption{Übersicht der Designvorschläge mit Herleitung und zugehörigen Playtest-Items}
    \label{tab:Designvorschlaege}
\end{table}

\begin{itemize}[leftmargin=1.5cm]
    \item \textcolor{red}{Tabelle prüfen und ggf. ergänzen; die Items Q9--Q11 sowie B5 und B10--B14 betreffen den gesamten Kampfabschnitt und keinen einzelnen Vorschlag.}
    \item \textcolor{red}{Optional: kurzer Satz, warum gerade diese acht Vorschläge ausgewählt wurden und was bewusst nicht umgesetzt wurde (z.\,B. Elemente aus den Analysen, die verworfen wurden, und weshalb).}
\end{itemize}

\subsection{Gegnerverhalten und Spieler-Gegner-Interaktion} \label{sec:UmsetzungGegner}

\subsubsection{Überblick über die Gegnertypen} \label{sec:Gegnertypen}
Im Unserem Prototyp gibt es zwei verschiedene Gegnerklassen, von welchen es verschiedene Varianten gibt. Die Gegner lassen sich in die Klassen Nah- und Fernkämpfer unterteilen. Von den Nahkämpfern gibt es zwei verschiedene Varianten und von den Fernkämpfern gibt es drei. Demnach gibt es insgesamt 5 verschiedene Gegnertzpen. [hier werde ich ein bild hinzufügen, von allen gegnern nebeneinander. bitte mach einen platzthalter rein damit ich es nicht vergesse]. Von links nach rechts sind die Codenamen für die gezeigten gegner: Solider, Heavy Solider, Archer, Heavy Archer, Chain Archer. Der Soldat ist der Der basic Nahkämpfer, der Auf den Spieler zurennt und schlägt. Der Heavy Solider ist grösser und stärker und hat mehr leben. er hat das gleiche verhalten wie der normale Solider. Die Archers bewegen sich alle nicht. Der normale Archer schiesst schiesst ein Projektil auf den Spieler. Der Heavy Archer macht das Selbe mit dem Unterschied, dass das Projektil langsamer ist aber die Richtung ändern kann und den Spieler verfolgt. Der Chain archer schiesst gleich 3 Pfeile auf den Spieler, die mehr als doppelt so schnell sind, wie die vom normalen Archer. Farblich sind die Nahkämpfer Rötlich und die Fernkämpfer Gelblich. Der Heavy Solider ist langsamer als der normale Solider. Die Reichweite der gegner ist gleich. Der Schaden, den die Gegner machen und wie viel leben sie haben ist unterschiedlich: Der Solider hat 130H und macht 15damage, der Heavy Soldier 200h/30d, der Pfeil vom Archer und die 3 vom Chain archer machen je 10 damage. Der ARcher hat 60 und der chain archer hat 70 leben. Der Heavy Archer hat 40 leben und macht 20 damage mit seinem Verfolgungspfeil. Der Speler selber hat 200 leben und seine attacken machen unterschiedlich viel schaden. Seine Heavy Combo macht 30, 30 40 schaden, seine Light Combo macht 20 20 30 40 schaden und seine air combo macht 20,20,20,20 schaden. Die schadenswerte und Lebenspunkte wurden so gewählt, dass wenn der spieler z.b. eine Heavy attack auf den Chain Archer macht, dass dieser mit dem letzten schlag der Combo stirbt, was wir gedacht haben, dass es sich besser anfühlt. oder das der Heavy Soldier nach genau 2 Heavy Combos stirbt. Wir haben aber auch darauf geachtet, dass z.b. Soliders gerade so nicht nach einer Light Combo sterben, damit beim spielen Variation entsteht und Spieler Combos abbrechen oder kombinieren müssen.
\begin{itemize}[leftmargin=1.5cm]
    \item \textcolor{red}{Alle Gegnertypen des Prototyps vorstellen, mit den Namen wie im Spiel (z.\,B. Soldier, Heavy Soldier, Archer, Heavy Archer, Chain Archer).}
    \item \textcolor{red}{Pro Typ: Rolle (Nah-/Fernkampf), Tempo, Reichweite, Schaden, Lebenspunkte, Angriffsmuster, Farbcode -- am besten als Tabelle analog zu Frage M1 der Analysen.}
    \item \textcolor{red}{Klar belegen, dass mindestens zwei unterscheidbare Gegnertypen umgesetzt sind (Vertrag).}
    \item \textcolor{red}{Screenshot je Gegnertyp (\enquote{Eigene Abbildung}).}
    \item \textcolor{red}{Gemeinsame Basisklasse und Verhaltensarchitektur nur kurz nennen (Rückbezug auf Abschnitt~\ref{sec:Kernmechanik}); Details mit konkretem Verweis auf die technische Dokumentation.}
\end{itemize}

\subsubsection{V1: Überlegenheit als Grundzustand} \label{sec:V1}

\paragraph{Vorschlag und Herleitung}
Der Spieler soll sich seinen Gegnern grundsätzlich überlegen fühlen, wobei dieser Zustand gezielt durchbrochen werden kann. Der Nutzen liegt im Spielgefühl: Aggression soll sich sicher und belohnt anfühlen, während Spannung erst dort entsteht, wo die Kontrolle kurzzeitig bedroht wird. Hergeleitet ist der Vorschlag aus dem Kontrast der beiden Referenzspiele, der dauerhaften Überlegenheit in \game{NieR: Automata} (Abschnitt~\ref{sec:ReferenzspielanalyseNier}) gegenüber dem prekären Gleichgewicht in \game{Ninja Gaiden} (Abschnitt~\ref{sec:ReferenzspielanalyseNinja}), sowie aus dem Prinzip des emotionsgeleiteten Designs (Abschnitt~\ref{sec:Pacing}).

\paragraph{Umsetzung im Prototyp}
Wie wir das in unserem spiel erreicht haben ist wie folgt: Zunächst ist der Spieler jedem einzelenn gegner in einem 1 gegen 1 haus hoch überlegen. Der Spieler macht meist doppelt bis vier mal so viel schaden wie die gegner und hat auch doppelt doppelt bis vier fach so viel leben mit der ausnahme vom heavy solider der mit 200 lebenspunkten genau gleich viel hat. zudem hat jeder gegner für sich ein simples moveset. heisst sie haben eine limitierte anzahl bewegungen die sie ausführen können. der Spieler hat deutlich mehr und kann z.b. schneller rennen, kann springen, was keiner der Gegner kann und kann sich sogar heilen. Man könnte sich deshalb denken, dass der Spieler keine Probleme haben wird, auch grosse Mengen von gegnern zu bewältigen. Das stimmt allerdings nicht ganz, da sich in unserem Spiel die gegner miteinander abstimmen können und so in der Gruppe eine grössere Gefahr darstellen. Zudem können einige Gegner nicht nur durch eine Combo vernichtet werden was grosse horden von Ihnen noch Gefährlicher macht. Des weiteren sind die allermeisten gegner ausser der Heavy Solider unterbrechbar, was sie im dierekten Nahkampf sehr sehr schwach macht. Gezielt unterbrochen ist die Überlegenheit des Spielers am ende einer combo. Das liegt daran, dass dieser eine Gewisse trägheit aufweist, welche bei nicht erfolgreichem killen des gegners dem gegner die möglichkeit gibt dem Spieler ein hit reinzudrücken. Das heilen des Spielers dauert auch eine gewisse zeit, heisst wenn es zu wenig Platzt gibt und keine möglichkeit sich zu verstecken, dann ist es sehr schwer zu heilen. Alternativ wie wir versucht haben die Encounters schwerer zu machen war, mit mehr leben für die gegner, mehr und längere Wellen. Wir haben uns aber gegen diese lösungen entschieden, da sie das spiel einfach nur sinnlos strecken und keinen mehrwert bieten. Sie verlangsamen dadurch den spielfluss, für welchen wir vorgesehen haben, dass dieser schnell und agressiv sein soll.

\begin{itemize}[leftmargin=1.5cm]
    \item \textcolor{red}{Wodurch fühlt sich der Spieler überlegen? (z.\,B. Schadenswerte, Unterbrechbarkeit der Gegner durch Treffer, Tempo Spieler vs. Gegner, Heilung)}
    \item \textcolor{red}{Wo wird die Überlegenheit gezielt durchbrochen? (z.\,B. Heavy Soldier, grosse Gruppen, enge Räume)}
    \item \textcolor{red}{Konkrete Werte/Parameter nennen, wo sinnvoll.}
    \item \textcolor{red}{Verworfene Alternativen und Begründung (Warum so? Welche Alternativen? Weshalb verworfen?).}
\end{itemize}

\paragraph{Überprüfbarkeit}
\begin{itemize}[leftmargin=1.5cm]
    \item \textcolor{red}{Szene/Datei/Klasse nennen bzw. auf das passende Kapitel der technischen Dokumentation verweisen.}
    \item \textcolor{red}{Playtest-Indikatoren: Q1, B1.}
\end{itemize}

\subsubsection{V2: Begrenzung gleichzeitiger Angreifer} \label{sec:V2}

\paragraph{Vorschlag und Herleitung}
Auch bei vielen Gegnern sollen nie alle gleichzeitig angreifen, sondern nur eine begrenzte Zahl. Der Nutzen betrifft Schwierigkeit und Spielfluss, da grosse Gruppen dadurch lesbar und fair bleiben, statt den Spieler zu überfordern. Beide Referenzspiele nutzen dieses Mittel, in \game{Ninja Gaiden} mit maximal etwa vier aktiven Gegnern, in \game{NieR: Automata} mit vermutlich höchstens zwei gleichzeitig angreifenden Nahkämpfern (Abschnitt~\ref{sec:Vergleich}), und es entspricht dem Pacing-Gedanken der Recherche (Abschnitt~\ref{sec:Pacing}).

\paragraph{Umsetzung im Prototyp}

Die Anzahl gegner, die gleichzeitig Spawnen ist wie folgt reguliert. Ein raum ist ein encounter. und dieser encounter hat mehrere wellen. Eine welle spawn erst nachdem alle Gegner der vorherigen gestorben sind. Wellen in unserem Prototyp können zwischen 10-30 Gegner enthalten. Wie jetzt die Anzahl gleichzeitig angreifender Gegner kontrolliert wird ist wie folgt. für jede welle kann man sich entscheiden, wie viele von der gesammten gegneranzahl gleichzeitig spawnen soll. Heisst z.b. wenn eine Welle 30 gegner hat und ich eingestellt habe dass nur 6 auf einmal im raum sein können dann sind nur 6 im raum. erst wenn einer stirbt, rutscht ein gegner nach. Des weiteren kann man vereinfacht gesagt entscheiden, wie agressiv eine welle ist und wie viele der gegner im angriffs modus oder im abstand halte modus sind. Die Grenzwerte und anzahl gegner und anzahl wellen haben wir je nach raumgrösse entschieden und angepasst, bis es uns gefallen hat. Wenn die gegner nicht angreifen, halten sie einen gewissen abstand zum Spieler und umkreisen ihn falls möglich. Archers die gar nicht laufen können bleiben natürlich einfach stehen.

\begin{itemize}[leftmargin=1.5cm]
    \item \textcolor{red}{Wie wird die Zahl gleichzeitiger Angreifer technisch begrenzt (z.\,B. zentrale Vergabe über einen Manager/Rollensystem)? Grenzwert nennen und begründen (Vergleich: \game{Ninja Gaiden} ca.\ vier, \game{NieR: Automata} ca.\ zwei).}
    \item \textcolor{red}{Was tun Gegner, die gerade nicht angreifen dürfen (warten, umkreisen, Position halten)?}
    \item \textcolor{red}{Gilt die Begrenzung auch für Fernkämpfer bzw. Projektile?}
\end{itemize}

\paragraph{Überprüfbarkeit}
\begin{itemize}[leftmargin=1.5cm]
    \item \textcolor{red}{Szene/Datei/Klasse bzw. Verweis auf die technische Dokumentation.}
    \item \textcolor{red}{Playtest-Indikatoren: Q2, B2.}
\end{itemize}

\subsubsection{V3: Ein taktisches Dilemma durch zwei Gegnerrollen} \label{sec:V3}

\paragraph{Vorschlag und Herleitung}
Der Kampfabschnitt soll Nah- und Fernkämpfer kombinieren, sodass der Spieler entscheiden muss, welche Bedrohung er zuerst ausschaltet. Der Nutzen liegt in der Schwierigkeit, die dadurch taktisch statt rein mechanisch (etwa durch mehr Lebenspunkte der Gegner) entsteht und schon kleine Gruppen anspruchsvoll macht. Grundlage ist, dass genau diese Kombination in beiden Referenzspielen den taktischen Druck trägt (Frage~D2 beider Analysen), sowie das Prinzip der Gegnerrollen im Encounter-Design (Abschnitt~\ref{sec:Encounter}).

\paragraph{Umsetzung im Prototyp}

Wie vorhin beschrieben haben wir 2 Nahkämpfer und 3 Fernkämpfer. Den Nahkämpfern kann man, wenn man will einfach wegrennen, da sie einen nur verfolgen, nicht springen könne und langsamer sind als der spieler. bei den vernkämpfern geht das nicht, weil sie einen nicht verfolgen sondern abschiessen. Deshalb stellen die Archers eine grosse gefahr dar. Der Grund wieso der spieler nicht einfach die Archers einen nach dem anderen killen kann besteht darin das aller meistens ein Nahkämpfer zwischen dem spieler und dem archer ist. das liegt daran, dass die Nahkämpfer sich bewegen und daruch dazwsichen kommen. deshalb ist der spieler oft gezwungen zuerst die Nahkämpfer auszuschalen das führt aber dazu dass die Archers dann zeit haben auf den Spieler zu schiessen. Was dann wiederum ein weiterer grund ist, weshalb der Spieler den archer killen möchte. Des weiteren ist für dieses Dillemma die platzierung der Archers wichtig, darauf gehen wir dann im nächsten kapitel ein aber hier ganz kurz: Die archers werden in unserem spiel oft hinter den Soliders und weit weg vom spieler gespawnt was dazu führt, dass der spielr sie nicht gleich töten kann.

\begin{itemize}[leftmargin=1.5cm]
    \item \textcolor{red}{Welche Gegnertypen übernehmen den Nahkampf, welche den Fernkampf (Verweis auf Abschnitt~\ref{sec:Gegnertypen})?}
    \item \textcolor{red}{Wie entsteht das Dilemma konkret? (z.\,B. Fernkämpfer hinter Nahkämpfern platziert, Nahkämpfer blockieren den Weg)}
    \item \textcolor{red}{Falls vorhanden: Rollensystem (Nahkämpfer/Flankierer/Fernkämpfer, vgl. Frage M2 zu \game{Ninja Gaiden}) und wie Rollen vergeben werden.}
\end{itemize}

\paragraph{Überprüfbarkeit}
\begin{itemize}[leftmargin=1.5cm]
    \item \textcolor{red}{Szene/Datei/Klasse bzw. Verweis auf die technische Dokumentation.}
    \item \textcolor{red}{Playtest-Indikatoren: Q3, B3.}
\end{itemize}

\subsubsection{V4: Klar telegrafierte Angriffe} \label{sec:V4}

\paragraph{Vorschlag und Herleitung}
Bevorstehende Gegnerangriffe sollen deutlich angekündigt werden, etwa durch sichtbare Vorzeichen. Der Nutzen betrifft Schwierigkeit und Spielgefühl, weil der Spieler auch im Gruppenkampf rechtzeitig reagieren kann und Treffer als eigene Fehler statt als Willkür erlebt. Vorbild ist das deutliche Telegrafieren in \game{NieR: Automata} (Frage~M2), theoretisch gestützt durch das Prinzip der Lesbarkeit (Abschnitt~\ref{sec:Encounter}).

\paragraph{Umsetzung im Prototyp}

alle unsere gegner haben eine visuelle Ankündigung vor dem angriff. das bezieht sich allerdings ausschliesslich auf die Animation und nicht auf SFX oder Parikel Effekte. Es wird immer eine kleine Animation abgespielt befor der gegner den spieler schlägt. Des weiteren werden die Angriffe von z.B. dem Soldier durch sein verhalten vorhergesagt. Das liegt daran, dass dieser Sobald er beim spieler ankommt und kurz inne hält, zuschlägt. Das ist ein sehr simples Muster, welches der Spieler beobachten kann. Wir haben uns dagegen entschieden, weiter audio-visuelle anzeichen neben den Animation zu machen weil zum einen uns die zeit gefehlt hat und zum anderen

\begin{itemize}[leftmargin=1.5cm]
    \item \textcolor{red}{Wie wird ein Angriff angekündigt? (Animation/Ausholbewegung, Leuchten/Farbe, Partikel, Sound; Zeitfenster zwischen Ankündigung und Treffer)}
    \item \textcolor{red}{Unterschiede zwischen den Gegnertypen (z.\,B. zielsuchende Geschosse des Heavy Archer).}
    \item \textcolor{red}{Screenshot eines Telegraphs (\enquote{Eigene Abbildung}).}
\end{itemize}

\paragraph{Überprüfbarkeit}
\begin{itemize}[leftmargin=1.5cm]
    \item \textcolor{red}{Szene/Datei/Klasse bzw. Verweis auf die technische Dokumentation.}
    \item \textcolor{red}{Playtest-Indikatoren: Q4, B4.}
\end{itemize}

\subsection{Levelaufbau und Gegnerplatzierung} \label{sec:UmsetzungLevel}

\subsubsection{Überblick über den Kampfabschnitt} \label{sec:Kampfabschnitt}
Unser Videospiel hat 12 verschiedene räume und jeder raum hat eine speziell konfigurierte anyahl welle und grösse von wellen. Jede welle hat einen gewissen teil fest platzierter und einen gewissen teil zufällige spawnender gegener. besonners das erste spwanen der gegner am anfang einer welle ist oft manuel von uns entschieden, um eine symetire in das spiel zu bekommen. Zudem sind die meisten archers random gespawn, da diese sich nicht bewegen können. Wenn gegner sterben, dann spwanen die weitere gegner nach. die reihenfolge in welcher gegner gespawnt werden ist schon im vorhinein festgelegt:
\begin{itemize}
    \item Raum 1: Welle 1 mit einem Solider um ihn vorzustellen und dann eine welle mit 5 solidern wo nur 2 auf einmal spawnen können
    \item Raum 2: Welle eins wie vorher auch nur mit einem Archer um ihn vorzustellen und danach auch wieder 5 wo nur 2 spawnen können.
    \item Raum 3: Das ist der raum wo zum ersten mal gegner in 3 wellen vermischt werden, es werden pro welle insgesamt etwa 10-15 gegner gespawnt und hier werden formationen und das dilemma vorgestellt.
    \item Raum 4: Hier wird der Heavy soldier zuerst einzeln vorgestellt auf engem raum. und dann gleich in kombindation mit dem Archer. Zudem wird hier zum ersten mal mit der box eine richtige platform vorgestellt, wo man sich vor den gegnern ausruhen kann. Dieser Raum hat insgesamt 3 Wellen. Die erste welle ist einfach ein troll als jump attack von einem einzigen solider.
    \item Raum 5: Dieser Raum ko
\end{itemize}
\begin{itemize}[leftmargin=1.5cm]
    \item \textcolor{red}{Gesamtaufbau: Raumabfolge, Anzahl Räume, Übergänge (z.\,B. Türen öffnen sich nach geräumtem Raum), Start und Ende.}
    \item \textcolor{red}{Molekül- oder Bubble-Diagramm des eigenen Kampfabschnitts (Methode aus Abschnitt~\ref{sec:Raumtypologien} und~\ref{sec:Prototyping} anwenden) sowie Draufsicht/Greybox-Screenshot.}
    \item \textcolor{red}{Vorgehen: Planung, Greybox (Trenchbroom), Iterationen; verworfene Varianten und Gründe.}
    \item \textcolor{red}{Geplante Spannungskurve über den gesamten Abschnitt (optional als Diagramm).}
    \item \textcolor{red}{Umsetzung in Godot nur kurz (Szenenaufbau, Spawner/Wellensystem); Details mit Verweis auf die technische Dokumentation.}
\end{itemize}

\subsubsection{V5: Lesbare Platzierung, dynamisches Verhalten} \label{sec:V5}

\paragraph{Vorschlag und Herleitung}
Gegner sollen zu Beginn lesbar und teils symmetrisch platziert sein, sich danach aber dynamisch verhalten, etwa den Spieler umzingeln. Der Nutzen betrifft Spielfluss und Spielgefühl: Die Anfangsaufstellung gibt Überblick, das anschliessende Verhalten hält den Kampf unter Druck und unvorhersehbar. Der Vorschlag verbindet die durchdachte, manuelle Platzierung aus \game{NieR: Automata} (Frage~A3) mit dem dynamischen Verhalten aus \game{Ninja Gaiden} (Frage~D3, Abschnitt~\ref{sec:Vergleich}) und stützt sich auf die Prinzipien der Lesbarkeit und der Gegner als raumbildende Elemente (Abschnitt~\ref{sec:Encounter}).

\paragraph{Umsetzung im Prototyp}
\begin{itemize}[leftmargin=1.5cm]
    \item \textcolor{red}{Startaufstellung der Gegner pro Raum: symmetrisch? vom Eingang aus sichtbar? Spawnpunkte manuell oder algorithmisch -- und warum.}
    \item \textcolor{red}{Wie verhalten sich die Gegner danach dynamisch (Umzingeln, Flankieren)?}
    \item \textcolor{red}{Abbildung mit eingezeichneten Spawnpunkten.}
\end{itemize}

\paragraph{Überprüfbarkeit}
\begin{itemize}[leftmargin=1.5cm]
    \item \textcolor{red}{Szene/Datei bzw. Verweis auf die technische Dokumentation.}
    \item \textcolor{red}{Playtest-Indikatoren: Q5, B6.}
\end{itemize}

\subsubsection{V6: Der Raum als Druckmittel} \label{sec:V6}

\paragraph{Vorschlag und Herleitung}
Die Grösse des Kampfraums soll gezielt zwischen Enge und Weite wechseln. Der Nutzen liegt in der Schwierigkeit, da ein enger Raum den Spieler zu methodischerem Spiel zwingt, während ein weiter Raum ihm taktische Distanz erlaubt. Vorbild ist die gezielte Raumverengung in \game{Ninja Gaiden} (Frage~M3), theoretisch untermauert durch die räumlichen Typologien Narrow und Prospect Space (Abschnitt~\ref{sec:Raumtypologien}).

\paragraph{Umsetzung im Prototyp}
\begin{itemize}[leftmargin=1.5cm]
    \item \textcolor{red}{Welche Räume sind eng, welche offen (Narrow vs.\ Prospect Space)? Ungefähre Grössen.}
    \item \textcolor{red}{Welche Gegnerkombinationen treten in welchem Raumtyp auf und warum?}
    \item \textcolor{red}{Rolle von Hindernissen und Strukturen (z.\,B. Boxen, Container, Brücke, Trap Room).}
\end{itemize}

\paragraph{Überprüfbarkeit}
\begin{itemize}[leftmargin=1.5cm]
    \item \textcolor{red}{Szene/Datei bzw. Verweis auf die technische Dokumentation.}
    \item \textcolor{red}{Playtest-Indikatoren: Q6, B7.}
\end{itemize}

\subsubsection{V7: Wellenstruktur mit Erholungsphasen} \label{sec:V7}

\paragraph{Vorschlag und Herleitung}
Der Kampfabschnitt soll in Wellen mit kontrollierter Eskalation und kurzen Erholungsphasen dazwischen aufgebaut sein. Der Nutzen betrifft den Spielfluss, indem sich Anspannung und Entlastung abwechseln und die Schwierigkeit planbar ansteigt. Beide Referenzspiele nutzen dieses Muster (Frage~D1 in \game{Ninja Gaiden}, Frage~A4 in \game{NieR: Automata}), das der Recherche als Rollercoaster-Pacing bekannt ist (Abschnitt~\ref{sec:Pacing}).

\paragraph{Umsetzung im Prototyp}
\begin{itemize}[leftmargin=1.5cm]
    \item \textcolor{red}{Wellenaufbau pro Raum: Anzahl Wellen, Gegner pro Welle, Auslöser der nächsten Welle.}
    \item \textcolor{red}{Wo liegen die Erholungsphasen (zwischen Wellen, zwischen Räumen)? Heilmöglichkeiten?}
    \item \textcolor{red}{Tabelle Raum $\times$ Welle $\times$ Gegnertypen.}
\end{itemize}

\paragraph{Überprüfbarkeit}
\begin{itemize}[leftmargin=1.5cm]
    \item \textcolor{red}{Szene/Datei bzw. Verweis auf die technische Dokumentation.}
    \item \textcolor{red}{Playtest-Indikatoren: Q7, B8.}
\end{itemize}

\subsubsection{V8: Gestaffelte Einführung neuer Gegner} \label{sec:V8}

\paragraph{Vorschlag und Herleitung}
Jeder neue Gegnertyp soll zuerst isoliert auftreten und erst danach mit anderen kombiniert werden. Der Nutzen betrifft Spielfluss und Schwierigkeit, weil der Spieler einen Gegner in Ruhe verstehen kann, bevor die Kombination ihn fordert. Grundlage ist das Scaffolding-Prinzip (Abschnitt~\ref{sec:Scaffolding}) sowie die entsprechende Struktur in \game{NieR: Automata} (Frage~D1).

\paragraph{Umsetzung im Prototyp}
\begin{itemize}[leftmargin=1.5cm]
    \item \textcolor{red}{Reihenfolge, in der die Gegnertypen eingeführt werden (Raum für Raum): wo erstmals isoliert, wo erstmals kombiniert?}
    \item \textcolor{red}{Bezug auf Scaffolding und Braiding (Abschnitt~\ref{sec:Scaffolding}).}
\end{itemize}

\paragraph{Überprüfbarkeit}
\begin{itemize}[leftmargin=1.5cm]
    \item \textcolor{red}{Szene/Datei bzw. Verweis auf die technische Dokumentation.}
    \item \textcolor{red}{Playtest-Indikatoren: Q8, B9.}
\end{itemize}

\end{document}
